% 《大学物理简明教程》第5版第9章解答题9.3.7—9.3.12
% 源书印刷/PDF第239—240页；核查杆运动几何和电流方向后独立推导。
% 9.3.9对源图几何给出计算值并保留电位正向的约定。

\subsection{解答题：转动导线、涡旋电场与互感}

\sourceproblem{9.3.7}{239--240}{239--240}
\problemtext 长$\ell$的直金属杆$ab$绕过$O$点的垂直轴匀速旋转，
角速度$\omega$；$aO=\ell/3$，$Ob=2\ell/3$，
匀强磁场$\bm B$与旋转轴同向。
求杆两端感应电势差的大小，比较哪端电势高。
\workings
取从$O$向$b$为正的杆坐标$x$，
$x_a=-\ell/3$、$x_b=2\ell/3$。
假设$\bm\omega$与$\bm B$同向，
各点局部动生电动势元为
$d\mathcal E=B\omega x\,dx$。
从$a$到$b$积分：
\[
\mathcal E_{a\to b}
=B\omega\int_{-\ell/3}^{2\ell/3}x\,dx
=\frac{B\omega}{2}
\left[\frac{4\ell^2}{9}-\frac{\ell^2}{9}\right]
=\boxed{\frac{B\omega\ell^2}{6}}.
\]
故
\[
\boxed{U_b-U_a=\frac{B\omega\ell^2}{6}}>0
\]
在该同向轴与电势约定下，\textbf{$b$端电势高}。
\crosscheck 若从图上观察角速度反转或把$\bm B$反向，
电势差有向符号必须反号；
大小为$B|\omega|\ell^2/6$，与杆所处的角位置无关。

\sourceproblem{9.3.8}{240}{240}
\problemtext 非洲某地地磁感应强度$B=4.0\times10^{-5}$ T，
磁场与\emph{水平面}的夹角为$60^\circ$。
将一根长$\ell=1.0$ m、沿东西方向水平放置的金属杆自由下落，
忽略阻力，求$t=3$ s时杆上动生电动势。
\workings
杆沿东西向，速度$\bm v$竖直向下，$v=gt$。
只有磁场的\emph{水平且南北向}分量使
$(\bm v\times\bm B)\cdot d\bm l$沿东西向非零。
已知地磁倾角$60^\circ$，水平场幅为
\[
\boxed{B_{\mathrm H}=B\cos60^\circ=\frac B2}.
\]
若水平场的方位角相对东西方向为$\beta$，
其南北向分量为$B_{\mathrm H}\sin\beta$，
故严谨的一般表达式为
\[
\boxed{|\mathcal E|=B\ell gt\,\cos60^\circ|\sin\beta|}.
\]
若采用通常教材默认的水平地磁场沿南北向
（$\beta=90^\circ$），则
\[
\boxed{|\mathcal E|=4.0\times10^{-5}(1.0)(9.8)(3)
\cdot\frac12
=5.88\times10^{-4}\,\mathrm V}.
\]
\crosscheck \textbf{地磁方位需二校：}源题只给与水平面的倾角，
未明说水平投影方位。在南北向水平分量的常规假设下得
$0.588\,\mathrm{mV}$；其他方位需乘$|\sin\beta|$。
此前将倾角的正弦误作水平分量会导致数值偏大，
本稿按与\emph{水平面}成角的定义使用\emph{余弦}修正。

\sourceproblem{9.3.9}{240}{240}
\problemtext 半径$R$的圆柱形区域内部有均匀、指向纸内的磁场，
$B(t)$随时间增大，圆外磁场为零。
一根长$2R$的水平金属杆$ac$横穿区域下部，
其前半段$ab$（长$R$）在圆内，
后半段$bc$（长$R$）在圆外；
原图半径$Ob$相对竖直方向倾斜$30^\circ$。
求两端感生电动势的大小和沿杆方向。
\workings
令圆心在$(0,0)$，磁场指向纸内，
取杆从$a$到$c$的水平方向为$x$正向。
源图中$Ob$与竖直方向夹角$30^\circ$，
所以水平弦与圆心的距离
$h=\sqrt3R/2$，并取
\[
a=(-R/2,-h),\qquad b=(R/2,-h),\qquad c=(3R/2,-h).
\]
记$\dot B=dB/dt>0$，表示纸内磁场\emph{大小}的增长率。
由轴对称法拉第环流，圆内的感生电场为
\[
\bm E_{\mathrm{in}}
=\frac{\dot B}{2}(-y\,\bm e_x+x\,\bm e_y).
\]
在$y=-h$处，$x$分量恒为$\dot Bh/2$，故
\[
\mathcal E_{a\to b}
=\frac{\dot Bh}{2}R
=\frac{\sqrt3}{4}\dot BR^2.
\]
圆外的感生电场切向分量
$E_\theta=\dot BR^2/(2r)$（沿逆时针），
沿水平线$x$方向的分量为
$E_x=\dot BR^2h/[2(x^2+h^2)]$。
因此
\[
\mathcal E_{b\to c}
=\frac{\dot BR^2}{2}
\left[
\arctan\frac{x}{h}
\right]_{R/2}^{3R/2}
=\frac{\pi}{12}\dot BR^2.
\]
合计
\[
\boxed{\mathcal E_{a\to c}
=\left(\frac{\sqrt3}{4}+\frac{\pi}{12}\right)R^2\dot B}.
\]
感生电场沿杆从\textbf{$a$指向$c$}；
以静电势差记忆时，要说明涡旋电场本身不可写成全局静电势梯度。
\crosscheck 本题必须分段计算\emph{圆内}与\emph{圆外}；
圆外尽管$B=0$，仍有非零涡旋电场。
数值系数基于图中$30^\circ$角确实相对\emph{竖直半径}的解释，
源图方位应二校。

\sourceproblem{9.3.10}{240}{240}
\problemtext 一根无限长直导线载交变电流
$I(t)=I_0\sin\omega t$。
在其右侧共面放一矩形线圈，线圈高$a$（与直导线平行）、
宽$b$，近边与直导线的距离为$d$。
求线圈中的感应电动势。
\workings
距直导线$x$处
$B(x,t)=\mu_0I(t)/(2\pi x)$，
以纸面向里为正磁通方向，有
\[
\Phi_B(t)=a\int_d^{d+b}B(x,t)\,dx
=\frac{\mu_0aI_0}{2\pi}
\ln\frac{d+b}{d}\sin\omega t.
\]
因此
\[
\boxed{\mathcal E(t)=
-\frac{\mu_0aI_0\omega}{2\pi}
\ln\frac{d+b}{d}\cos\omega t}.
\]
大小峰值
$\boxed{\mathcal E_{\max}
=\mu_0aI_0\omega\ln[(d+b)/d]/(2\pi)}$。
\crosscheck 频率升高会线性提高同一$I_0$条件下的电动势峰值；
固定导线的磁场具有$1/x$的径向分布，不能取为均匀场。

\sourceproblem{9.3.11}{240}{240}
\problemtext 图9.3.11中无限长导线穿过一矩形线圈所在的平面，
长导线电流$I(t)=I_0\sin\omega t$；
矩形高度$a$，位于长导线左侧的宽度$a/2$、
右侧宽度$3a/4$。求互感系数及互感电动势。
\workings
长导线两侧磁场指向纸面的方向相反。
用线圈法向统一定向，两个区域磁通量的有向贡献相减。
在导线离线圈高度方向上取单位长，
可把两侧径向积分的公共小半径截止量$\epsilon$
逐项消去：
\[
\Phi_B=\frac{\mu_0Ia}{2\pi}
\left[
\ln\frac{3a/4}{\epsilon}
-\ln\frac{a/2}{\epsilon}
\right]
=\frac{\mu_0Ia}{2\pi}\ln\frac32.
\]
故互感系数大小
\[
\boxed{M=\frac{\mu_0a}{2\pi}\ln\frac32},
\]
互感电动势为
\[
\boxed{\mathcal E(t)=-M\frac{dI}{dt}
=-\frac{\mu_0aI_0\omega}{2\pi}
\ln\frac32\cos\omega t}.
\]
\crosscheck 当左右宽度相等时两侧磁通完全抵消，互感为零；
此处正是两侧宽度不同才出现非零净磁链。
实际长导线有有限半径，以上推导中两边同一截止尺度相消。

\sourceproblem{9.3.12}{240}{240}
\problemtext 一对相互平行的长直导线，其横截面圆半径同为$a$，
两导线中心相距$d$，组成一往返电路。
忽略导线\emph{内部}磁通量，证明每段长度$\ell$的自感为
$L=(\mu_0\ell/\pi)\ln[(d-a)/a]$。
\workings
两导线载电流大小$I$、方向相反，
对两导线间跨径向的狭长曲面，
两导线磁场沿法向\emph{同向}。
从靠近第一根导线表面$r=a$积分到
靠近第二根表面$r=d-a$：
\[
\Phi_B=\frac{\mu_0 I\ell}{2\pi}
\int_a^{d-a}\left(\frac1r+\frac1{d-r}\right)dr
=\frac{\mu_0I\ell}{\pi}\ln\frac{d-a}{a}.
\]
因此
\[
\boxed{L=\frac{\Phi_B}{I}
=\frac{\mu_0\ell}{\pi}\ln\frac{d-a}{a}}.
\]
\crosscheck 需要$d>2a$使两圆柱互不相交；
该式是外部磁链的几何近似，
不包含两导线内部的电流分布所带来的内部自感。
